\documentclass[10pt,a4paper]{amsart}
\usepackage[margin=19mm]{geometry}
\usepackage{amsmath,amssymb,mathtools}
\usepackage[scr=rsfs,cal=euler]{mathalfa}
\usepackage{xfrac}
\usepackage[unicode,hidelinks,bookmarks=false]{hyperref}
\newcommand{\ie}{i.e.\ }
\newcommand{\cf}{cf.\ }
\newcommand{\resp}{respectively\ }
\newcommand{\Subsection}{\subsection}
\providecommand{\nobreakdash}{\nobreak\hbox{-}}
\setlength{\emergencystretch}{2em}
\multlinegap0pt
\usepackage{amssymb}
\usepackage{enumerate}
\usepackage{csquotes}
\newtheorem{lem}{Lemma}
\newtheorem{prop}{Proposition}
\newcommand{\Irr}{{\operatorname{Irr}}}
\title{Zeros of $S$-characters}
\author{Thomas Breuer, Michael Joswig, Gunter Malle}
\date{Source: arXiv:2408.16785v3 (CC BY 4.0)}
\begin{document}
\maketitle

\noindent\emph{Source and licence.} Zeros of $S$-characters. Version arXiv:2408.16785v3 (CC BY 4.0), distributed by arXiv under the Creative Commons Attribution 4.0 International licence, http://creativecommons.org/licenses/by/4.0/, as recorded in the arXiv licence metadata for this exact version and shown on the abstract page https://arxiv.org/abs/2408.16785v3. Full licence notice: \texttt{licenses/arxiv-2408.16785-CC-BY-4.0.txt}.\par\smallskip

\noindent\textit{Editorial scope.} Counted unit: the article's prop prop:solvnotord (source0.tex lines 444--447). The article's complete original proof of it is delivered with it (source0.tex lines 449--467, one \texttt{proof} environment of 19 lines). No mathematics is written, repaired or re-derived: the statement and the proof are the article's own lines, in the article's own notation, and no retained line is substituted or re-flowed.  Retained uncounted as the source-local context the proof needs: lines 405--409.  Nothing was omitted from any retained span, so the package carries no omission record.  No retained line contains a citation, so the wrapper carries no bibliography and no bibliography entry.  No retained line contains a figure environment, an image-inclusion command or drawing code, so no image is delivered and the package declares no asset dependency.  Editorial formatting only, in addition: an AMS \texttt{amsart} preamble that restates the article's own declarations and macros used by the retained lines, namely the environments \texttt{lem}, \texttt{prop} on separate counters, together with the standard packages those lines need.  The retained environments print in this document as Lemma 1, Proposition 1 in order of appearance, whereas the article prints the same items with the numbering of its own full document.  The complete native source of the article as distributed in the arXiv e-print for this exact version is bundled unchanged as \texttt{sources/164481/source0.tex}.
\begin{lem}   \label{lem:proj}
 Let $\chi\in\Irr(G)$. Then either $\chi^F\in \Irr(F)$ or $\chi^F$ is zero, and
 $$(\chi^F)_G=\begin{cases} \chi& \text{if $N\subseteq \ker(\chi)$},\\
  0& \text{otherwise}.\end{cases}$$
\end{lem}

\begin{prop}\label{prop:solvnotord}
 Each non-trivial ordinary $S$-character of a finite solvable group vanishes on
 some element of prime power order.
\end{prop}

\begin{proof}
Let $G$ be a counterexample of minimal order, and let $\psi$ be a non-trivial
ordinary $S$-character of $G$ that is nonzero on all elements of prime power
order. Since $\psi$ vanishes at some element, $G$ is not elementary abelian.

Let $N$ be a non-trivial elementary abelian (and hence proper) normal subgroup
of~$G$, and set $F = G/N$. By Lemma~\ref{lem:proj} and the subsequent remark,
$\psi^F$ is an $S$-character of~$F$ that is nonzero on all elements of prime
power order. By the minimality of~$G$, $\psi^F$ is the trivial character.

This means that $\psi$ consists of the trivial character plus irreducible
constituents $\chi$ with the property $N \not\subseteq \ker(\chi)$.
By the theorem of Clifford, the restriction of any such $\chi$ to $N$ does not
have a trivial constituent, which implies that the restriction of $\psi$ to $N$
is an $S$-character of $N$ that is nonzero on all elements of prime power order.
As above, we conclude that this restriction is the trivial character.
Since $\psi$ is an ordinary character with a trivial constituent,
$\psi$ itself is trivial.
\end{proof}

\end{document}
